% GENERATED FILE. Edit canonical lessons, then run npm run generate:books.
% canonical-source-hash: fnv1a64:83d20f53eb9fd4ac
% canonical-lessons: UR-C514-un, UR-C514-pirona, UR-C514-karhna, UR-C514-angushtana, UR-C514-jhalar

\chapter{Needle and Pattern}
\label{ch:ur-needle-and-pattern}
\hlchaptermodality{514}

\begin{chapteropening}
\textbf{By the end of this chapter:} \emph{I can talk about embroidery and handwork.}

Complete the last lesson of chapter 514: I can talk about embroidery and handwork.
\end{chapteropening}

\section[\texorpdfstring{\ur{اون} (ūn)}{اون (ūn)}]{\ur{اون} (ūn) — wool}

\label{lesson:UR-C514-un}



\begin{quote}
Before the new one: say the Urdu for a picnic, then the Urdu for an album.
\end{quote}

\subsection*{You'll want to know: \ur{اون}}
\textbf{\ur{اون}} — \emph{ūn} — \textquotedblleft{}wool\textquotedblright{}.

\begin{quote}
\ur{سردیوں} \ur{میں} \ur{ہم} \ur{اون} \ur{کے} \ur{سویٹر} \ur{پہنتے} \ur{ہیں۔}
\end{quote}

\begin{quote}
\emph{sardiyoṅ meṅ ham ūn ke sveṭar pahante haiṅ.}
\end{quote}

\begin{quote}
In winter we wear woollen jumpers.
\end{quote}

Feminine.

\begin{cousinweb}
\textbf{\ur{اون}} (\emph{ūn}) is inherited from Sanskrit \emph{ūrṇā}, \textquotedblleft{}wool\textquotedblright{} (CDIAL 2424).
\end{cousinweb}

\begin{grammarlens}[title={What you've built}]
One more word for describing what you lived and felt.
\end{grammarlens}

\subsection*{Your turn}
\begin{itemize}
\raggedright
  \item \emph{Say it:} \emph{ūn}
  \item \emph{Say it:} \emph{ūn}, once more
  \item \emph{Say it:} \emph{sardiyoṅ meṅ ham ūn ke sveṭar pahante haiṅ.} — In winter we wear woollen jumpers.
  \item \emph{Recall:} say the Urdu for a picnic, then the Urdu for an album, then say \emph{ūn} again
\end{itemize}

\begin{tcolorbox}[breakable,colback=teal!4,colframe=teal!35!black,title={Before you move on}]
What does \ur{اون} mean? (\textquotedblleft{}Wool\textquotedblright{}.) Say it once more.
\end{tcolorbox}
\section[\texorpdfstring{\ur{پرونا} (pironā)}{پرونا (pironā)}]{\ur{پرونا} (pironā) — to thread (a needle); to string (beads)}

\label{lesson:UR-C514-pirona}



\begin{quote}
Before the new one: say the Urdu for an album, then the Urdu for wool.
\end{quote}

\subsection*{You'll want to know: \ur{پرونا}}
\textbf{\ur{پرونا}} — \emph{pironā} — \textquotedblleft{}to thread (a needle); to string (beads)\textquotedblright{}.

\begin{quote}
\ur{نانی} \ur{نے} \ur{سوئی} \ur{میں} \ur{دھاگا} \ur{پرویا۔}
\end{quote}

\begin{quote}
\emph{nānī ne sūī meṅ dhāgā piroyā.}
\end{quote}

\begin{quote}
Grandmother threaded the needle.
\end{quote}

Past, with \ur{نے} (he or she …): \ur{اس} \ur{نے} \ur{سوئی} \ur{میں} \ur{دھاگا} \ur{پرویا}, \emph{us ne sūī meṅ dhāgā piroyā}.

\begin{cousinweb}
\textbf{\ur{پرونا}} (\emph{pironā}) is inherited: CDIAL lists the Hindi cousin \emph{pironā}, \textquotedblleft{}to string pearls, to thread a needle\textquotedblright{}, under Sanskrit \emph{parivayati} (CDIAL 7869).
\end{cousinweb}

\begin{grammarlens}[title={What you've built}]
One more word for describing what you lived and felt.
\end{grammarlens}

\subsection*{Your turn}
\begin{itemize}
\raggedright
  \item \emph{Say it:} \emph{pironā}
  \item \emph{Say it:} \emph{pironā}, once more
  \item \emph{Say it:} \emph{nānī ne sūī meṅ dhāgā piroyā.} — Grandmother threaded the needle.
  \item \emph{Recall:} say the Urdu for an album, then the Urdu for wool, then say \emph{pironā} again
\end{itemize}

\begin{tcolorbox}[breakable,colback=teal!4,colframe=teal!35!black,title={Before you move on}]
What does \ur{پرونا} mean? (\textquotedblleft{}To thread (a needle); to string (beads)\textquotedblright{}.) Say it once more.
\end{tcolorbox}
\section[\texorpdfstring{\ur{کاڑھنا} (kāṛhnā)}{کاڑھنا (kāṛhnā)}]{\ur{کاڑھنا} (kāṛhnā) — to embroider}

\label{lesson:UR-C514-karhna}



\begin{quote}
Before the new one: say the Urdu for wool, then the Urdu for to thread (a needle).
\end{quote}

\subsection*{You'll want to know: \ur{کاڑھنا}}
\textbf{\ur{کاڑھنا}} — \emph{kāṛhnā} — \textquotedblleft{}to embroider\textquotedblright{}.

\begin{quote}
\ur{امی} \ur{نے} \ur{دوپٹے} \ur{پر} \ur{پھول} \ur{کاڑھے۔}
\end{quote}

\begin{quote}
\emph{ammī ne dupaṭṭe par phūl kāṛhe.}
\end{quote}

\begin{quote}
Mum embroidered flowers on the scarf.
\end{quote}

Past, with \ur{نے} (he or she …): \ur{اس} \ur{نے} \ur{تکیے} \ur{پر} \ur{پھول} \ur{کاڑھا}, \emph{us ne takiye par phūl kāṛhā}.

\begin{cousinweb}
\textbf{\ur{کاڑھنا}} (\emph{kāṛhnā}) is inherited: CDIAL lists the Hindi cousin \emph{kāṛhnā}, \textquotedblleft{}to pull out\textquotedblright{}, under \emph{*kaḍḍhati}, \textquotedblleft{}pulls, draws\textquotedblright{} (CDIAL 2660); the sense \textquotedblleft{}to embroider\textquotedblright{}, drawing a thread through cloth, is a later one that CDIAL does not record.
\end{cousinweb}

\begin{grammarlens}[title={What you've built}]
One more word for describing what you lived and felt.
\end{grammarlens}

\subsection*{Your turn}
\begin{itemize}
\raggedright
  \item \emph{Say it:} \emph{kāṛhnā}
  \item \emph{Say it:} \emph{kāṛhnā}, once more
  \item \emph{Say it:} \emph{ammī ne dupaṭṭe par phūl kāṛhe.} — Mum embroidered flowers on the scarf.
  \item \emph{Recall:} say the Urdu for wool, then the Urdu for to thread (a needle), then say \emph{kāṛhnā} again
\end{itemize}

\begin{tcolorbox}[breakable,colback=teal!4,colframe=teal!35!black,title={Before you move on}]
What does \ur{کاڑھنا} mean? (\textquotedblleft{}To embroider\textquotedblright{}.) Say it once more.
\end{tcolorbox}
\section[\texorpdfstring{\ur{انگشتانہ} (angushtānā)}{انگشتانہ (angushtānā)}]{\ur{انگشتانہ} (angushtānā) — a thimble}

\label{lesson:UR-C514-angushtana}



\begin{quote}
Before the new one: say the Urdu for to thread (a needle), then the Urdu for to embroider.
\end{quote}

\subsection*{You'll want to know: \ur{انگشتانہ}}
\textbf{\ur{انگشتانہ}} — \emph{angushtānā} — \textquotedblleft{}a thimble\textquotedblright{}.

\begin{quote}
\ur{نانی} \ur{کی} \ur{انگلی} \ur{میں} \ur{ایک} \ur{پرانا} \ur{انگشتانہ} \ur{ہے۔}
\end{quote}

\begin{quote}
\emph{nānī kī unglī meṅ ek purānā angushtānā hai.}
\end{quote}

\begin{quote}
Grandmother has an old thimble on her finger.
\end{quote}

Masculine.

\begin{cousinweb}
\textbf{\ur{انگشتانہ}} (\emph{angushtānā}) is usually traced to Persian \emph{angushtāna}, \textquotedblleft{}a thimble\textquotedblright{}, built on \emph{angusht}, \textquotedblleft{}a finger\textquotedblright{}.
\end{cousinweb}

\begin{grammarlens}[title={What you've built}]
One more word for describing what you lived and felt.
\end{grammarlens}

\subsection*{Your turn}
\begin{itemize}
\raggedright
  \item \emph{Say it:} \emph{angushtānā}
  \item \emph{Say it:} \emph{angushtānā}, once more
  \item \emph{Say it:} \emph{nānī kī unglī meṅ ek purānā angushtānā hai.} — Grandmother has an old thimble on her finger.
  \item \emph{Recall:} say the Urdu for to thread (a needle), then the Urdu for to embroider, then say \emph{angushtānā} again
\end{itemize}

\begin{tcolorbox}[breakable,colback=teal!4,colframe=teal!35!black,title={Before you move on}]
What does \ur{انگشتانہ} mean? (\textquotedblleft{}A thimble\textquotedblright{}.) Say it once more.
\end{tcolorbox}
\section[\texorpdfstring{\ur{جھالر} (jhālar)}{جھالر (jhālar)}]{\ur{جھالر} (jhālar) — a fringe (on cloth)}

\label{lesson:UR-C514-jhalar}



\begin{quote}
Before the new one: say the Urdu for to embroider, then the Urdu for a thimble.
\end{quote}

\subsection*{You'll want to know: \ur{جھالر}}
\textbf{\ur{جھالر}} — \emph{jhālar} — \textquotedblleft{}a fringe (on cloth)\textquotedblright{}.

\begin{quote}
\ur{دلہن} \ur{کے} \ur{دوپٹے} \ur{پر} \ur{سنہری} \ur{جھالر} \ur{ہے۔}
\end{quote}

\begin{quote}
\emph{dulhan ke dupaṭṭe par sunahrī jhālar hai.}
\end{quote}

\begin{quote}
The bride's scarf has a golden fringe.
\end{quote}

Feminine.

\begin{cousinweb}
\textbf{\ur{جھالر}} (\emph{jhālar}) is inherited: CDIAL lists the Hindi cousin \emph{jhālar}, \textquotedblleft{}a fringe\textquotedblright{}, under Sanskrit \emph{jhallarī}, \textquotedblleft{}a curl\textquotedblright{} (CDIAL 5356).
\end{cousinweb}

\begin{grammarlens}[title={What you've built}]
One more word for describing what you lived and felt.
\end{grammarlens}

\subsection*{Your turn}
\begin{itemize}
\raggedright
  \item \emph{Say it:} \emph{jhālar}
  \item \emph{Say it:} \emph{jhālar}, once more
  \item \emph{Say it:} \emph{dulhan ke dupaṭṭe par sunahrī jhālar hai.} — The bride's scarf has a golden fringe.
  \item \emph{Recall:} say the Urdu for to embroider, then the Urdu for a thimble, then say \emph{jhālar} again
\end{itemize}

\begin{tcolorbox}[breakable,colback=teal!4,colframe=teal!35!black,title={Before you move on}]
What does \ur{جھالر} mean? (\textquotedblleft{}A fringe (on cloth)\textquotedblright{}.) Say it once more.
\end{tcolorbox}
